%======================================================================
%  Community Security Alert System (CSAS)
%  Systems Analysis and Design -- Workshop No. 2 (Revised Version)
%
%  Universidad Distrital Francisco Jose de Caldas
%  Computer Engineering Program -- Semester 2026-I
%
%  REVISION NOTE
%  This version addresses all feedback points from the first submission:
%   (1) System architecture (org + tech) added alongside software arch
%   (2) Concept of Operations (CONOPS) -- human-process view
%   (3) Design Decision Records (DDR-01..03) with alternatives & rationale
%   (4) Zone/time dispatch model formalised with weights, thresholds,
%       and three worked operational scenarios
%   (5) Technical-report format replacing academic-paper format
%   (6) Every design decision explicitly traced to W1 empirical evidence
%======================================================================

\documentclass[12pt,a4paper]{article}

%------ Packages -------------------------------------------------------
\usepackage[margin=2.5cm,top=2.8cm,bottom=2.8cm]{geometry}
\usepackage[utf8]{inputenc}
\usepackage[T1]{fontenc}
\usepackage{lmodern}
\usepackage{microtype}
\usepackage{graphicx}
\usepackage{booktabs}
\usepackage{array}
\usepackage{tabularx}
\usepackage{multirow}
\usepackage{longtable}
\usepackage{enumitem}
\usepackage{amsmath}
\usepackage{xcolor}
\usepackage{fancyhdr}
\usepackage{titlesec}
\usepackage{caption}
\usepackage[hidelinks]{hyperref}

\usepackage{tikz}
\usetikzlibrary{shapes.geometric,arrows.meta,positioning,calc,
                fit,matrix,decorations.pathmorphing}

%------ Colors ---------------------------------------------------------
\definecolor{udblue}{HTML}{003087}
\definecolor{udgold}{HTML}{C8962E}
\definecolor{soft}{HTML}{F5F5F5}
\definecolor{rowA}{HTML}{EAF1FB}
\definecolor{rowB}{HTML}{FFFFFF}
\definecolor{critcol}{HTML}{B53737}
\definecolor{highcol}{HTML}{C07020}
\definecolor{stdcol}{HTML}{1B7F3A}

%------ Page style -----------------------------------------------------
\pagestyle{fancy}
\fancyhf{}
\fancyhead[L]{\small\textcolor{udblue}{CSAS --- Systems Design Report --- Workshop No.\ 2 (Revised)}}
\fancyhead[R]{\small\textcolor{udblue}{Semester 2026-I}}
\fancyfoot[C]{\small\thepage}
\renewcommand{\headrulewidth}{0.4pt}

%------ Section formatting ---------------------------------------------
\titleformat{\section}[block]
  {\large\bfseries\color{udblue}}{\thesection.}{0.5em}{}[\vspace{-4pt}\rule{\linewidth}{0.4pt}]
\titleformat{\subsection}[runin]
  {\normalsize\bfseries}{\thesubsection.}{0.5em}{}[\ ---\ ]
\titleformat{\subsubsection}[runin]
  {\normalsize\itshape}{\thesubsubsection.}{0.5em}{}[\ ---\ ]

%------ Misc -----------------------------------------------------------
\renewcommand{\arraystretch}{1.2}
\setlist[enumerate]{itemsep=2pt,parsep=0pt}
\setlist[itemize]{itemsep=2pt,parsep=0pt}
\captionsetup{font=small,labelfont=bf}

%======================================================================
\begin{document}

%------ Title page -----------------------------------------------------
\begin{titlepage}
\centering
\vspace*{1.5cm}
{\color{udblue}\rule{\linewidth}{1.5pt}}\vspace{6pt}
{\color{udgold}\rule{\linewidth}{4pt}}\vspace{16pt}

{\LARGE\bfseries\color{udblue}
Community Security Alert System (CSAS)}\\\vspace{8pt}
{\large\bfseries Systems Design Report}\\\vspace{4pt}
{\large Workshop No.\ 2 --- Revised Submission}\\\vspace{16pt}

{\color{udgold}\rule{\linewidth}{4pt}}\vspace{6pt}
{\color{udblue}\rule{\linewidth}{1.5pt}}\vspace{24pt}

\begin{tabular}{ll}
\textbf{Authors} & Felipe José Garzón Herrera \hfill \textit{20251020132}\\
                 & Juan Esteban Quintero Gordillo \hfill \textit{20251020137}\\
                 & Henry Samuel Garrido Medina \hfill \textit{20251020125}\\
                 & Gabriel Mateo Cusba Marín \hfill \textit{20251020128}\\[6pt]
\textbf{Course}  & Systems Analysis and Design\\
\textbf{Professor} & Eng.\ Carlos Andrés Sierra, M.Sc.\\
\textbf{Program} & Computer Engineering\\
\textbf{Institution} & Universidad Distrital Francisco José de Caldas\\
\textbf{Semester} & 2026-I\\
\end{tabular}

\vfill
{\small\color{gray}
Document type: Systems Engineering Technical Report.
This revision addresses professor feedback on the first submission
and converts the document from academic-paper format to systems
engineering report format per workshop guidelines.}
\end{titlepage}

%------ Table of contents ----------------------------------------------
\tableofcontents
\newpage

%======================================================================
\section{Introduction and Design Rationale}

This document constitutes the systems design report for the Community
Security Alert System (CSAS), a socio-technical intervention on the
campus safety system of Universidad Distrital Francisco José de Caldas.
The design is grounded exclusively in the empirical findings of
Workshop~No.~1, which characterised campus safety as a
system bounded not by the willingness of its members to act but by the
\emph{latency}, \emph{reliability}, and \emph{traceability} of the
information flow connecting incident, awareness, and response.

Table~\ref{tab:gaps} restates the operational gaps from Workshop~No.~1
that drive every design decision in this document.

\begin{table}[h]
\caption{Operational Gap Analysis --- W1 Baseline and W2 Design Targets}
\label{tab:gaps}
\centering
\begin{tabularx}{\textwidth}{lXXX}
\toprule
\textbf{Dimension} & \textbf{Baseline (W1)} & \textbf{Design target} & \textbf{Gap} \\
\midrule
Incident-to-response time & $\approx$11.4 min & $<$5 min & $-$6.4 min \\
Incident record completeness & $\approx$41\% of events & $\geq$85\% & $+$44 pp \\
Community channel adoption & $\approx$12\% & $\geq$60\% & $+$48 pp \\
Alert delivery success & $\approx$87\% (phone/SMS) & $\geq$99.9\% & $+$12.9 pp \\
\bottomrule
\end{tabularx}
\end{table}

The document is organised as a systems engineering report, not an
academic paper: it prioritises traceable design decisions, operational
workflows, system architecture (people + processes + technology), and
concrete usage scenarios over abstract description.

%======================================================================
\section{Requirements Analysis}

\subsection{Functional requirements.}
Table~\ref{tab:fr} lists the ten functional requirements.
Each is explicitly traced to the W1 empirical finding that
motivates it, operationalising the principle that requirements must
emerge from system analysis.

\begin{table}[h!]
\caption{Functional Requirements with W1 Traceability}
\label{tab:fr}
\centering
\begin{tabularx}{\textwidth}{lXX}
\toprule
\textbf{ID} & \textbf{Requirement} & \textbf{W1 evidence that motivates it} \\
\midrule
FR-01 & Users shall report incidents by type, GPS location, and description in $\leq$3 interactions. & High-friction reporting identified as root cause of under-reporting (interview + observation). \\
FR-02 & Real-time geofenced alerts shall broadcast to users within 500\,m of a verified incident. & Spatial concentration of incidents in specific zones (observation: 2/2 suspicious events in parking lot at night). \\
FR-03 & An AI plausibility scorer shall filter implausible reports before human review. & Need to mitigate alert fatigue (Balancing Loop B1, W1 CLD). \\
FR-04 & Security staff shall have a web dashboard to review, escalate, and resolve reports. & Staff interviews: no centralised incident management tool exists. \\
FR-05 & Community members shall confirm or flag reports submitted by others. & 100\% of incident-experienced survey respondents willing to participate (W1 survey, $n=20$). \\
FR-06 & The mobile client shall capture GPS coordinates automatically on report submission. & Observer note: verbal descriptions insufficient for zone-level routing; exact location needed. \\
FR-07 & Every incident shall be persisted in an auditable log with timestamps and outcomes. & Institutional document analysis: formal records capture only $\approx$41\% of actual events. \\
FR-08 & A one-touch SOS button shall trigger immediate dispatch bypassing verification. & Interview: average $\approx$11.4-min latency incompatible with life-safety emergencies. \\
FR-09 & The analytics engine shall generate zone-level security heatmaps for university management. & Benchmarking: Policía Nacional / SIURE provide no real-time analytical layer. \\
FR-10 & Users shall attach image or video evidence to reports. & Observation: suspicious events went unrecorded by other witnesses; media adds verifiability. \\
\bottomrule
\end{tabularx}
\end{table}

\subsection{Non-functional requirements.}
\begin{itemize}
\item \textbf{NFR-01 (Latency):} End-to-end latency from submission to alert broadcast $<$30\,s (95th percentile).
\item \textbf{NFR-02 (Reliability):} Notification delivery $\geq$99.9\% across all channels.
\item \textbf{NFR-03 (Availability):} System uptime 99.9\%; guaranteed availability 18:00--22:00 (highest-risk window per W1 observation).
\item \textbf{NFR-04 (Privacy):} Full compliance with Ley 1581 de 2012; TLS~1.3 for all data in transit.
\item \textbf{NFR-05 (Scalability):} API gateway sustains a 300\% concurrent-request surge without performance degradation.
\item \textbf{NFR-06 (Usability):} Reporting workflow completable in $\leq$3 interactions; onboarding in $\leq$2 minutes.
\item \textbf{NFR-07 (Maintainability):} Any service component shall be independently updatable.
\item \textbf{NFR-08 (Integrity):} No alert dispatches as verified without $\geq$2 community confirmations \emph{or} 1 administrative approval.
\end{itemize}

%======================================================================
\section{Concept of Operations (CONOPS)}
\label{sec:conops}

Before specifying architecture, it is essential to describe \emph{how
the system operates from a human perspective}: who does what, under
what conditions, and with what information. This section fulfils the
systems engineering obligation to model the interaction of people,
processes, and technology before defining the technology itself.

\subsection{Operational roles.}
Table~\ref{tab:roles} defines the four roles that execute the
operational workflows described below.

\begin{table}[h]
\caption{Operational Roles in the CSAS Environment}
\label{tab:roles}
\centering
\begin{tabularx}{\textwidth}{lXX}
\toprule
\textbf{Role} & \textbf{Description} & \textbf{Primary system interaction} \\
\midrule
Community member & Any student, professor, or staff member with a registered account & Submits reports, receives geofenced alerts, confirms or flags reports from others \\
On-duty security guard & Campus security team member scheduled for a shift & Receives high-priority and emergency alerts; approves or rejects pending reports; coordinates ground response \\
Security supervisor & Duty officer responsible for escalation decisions & Manages the administrative dashboard; escalates verified incidents to local authorities \\
System administrator & IT staff maintaining the platform & Manages accounts, monitors health, accesses audit logs under controlled protocol \\
\bottomrule
\end{tabularx}
\end{table}

\subsection{Normal incident workflow.}
Fig.~\ref{fig:conops} presents the CONOPS swim-lane diagram covering
the standard incident lifecycle from report submission to resolution.
The diagram makes the interaction between human processes (left lanes)
and technological components (right lanes) explicit.

\begin{figure*}[ht!]
\centering
\begin{tikzpicture}[
  font=\footnotesize,
  every node/.style={align=center},
  swimlane/.style={draw=udblue!40,fill=udblue!5,
                   minimum width=2.6cm,minimum height=0.6cm,
                   rounded corners=2pt},
  tech/.style={draw=udgold!60,fill=udgold!10,
               minimum width=2.6cm,minimum height=0.6cm,
               rounded corners=2pt},
  arr/.style={-{Latex[length=2mm]},thick},
  cond/.style={diamond,draw=black!70,fill=soft,
               inner sep=1pt,aspect=2,font=\scriptsize}
]

% --- Header labels ---
\node[swimlane,minimum width=2.6cm] (h1) at (0,0)     {\textbf{Community\\Member}};
\node[swimlane,minimum width=2.6cm] (h2) at (3.4,0)   {\textbf{Mobile App /\\Web Portal}};
\node[tech,minimum width=2.6cm]     (h3) at (6.8,0)   {\textbf{Backend\\Services}};
\node[swimlane,minimum width=2.6cm] (h4) at (10.2,0)  {\textbf{Security\\Guard}};
\node[swimlane,minimum width=2.6cm] (h5) at (13.6,0)  {\textbf{Local\\Authorities}};

% --- Background lanes (drawn first so nodes appear on top) ---
\fill[udblue!4]  (-1.3,-0.35) rectangle (1.3,-10.2);
\fill[udgold!6]  (2.1,-0.35)  rectangle (4.7,-10.2);
\fill[udgold!10] (5.5,-0.35)  rectangle (8.1,-10.2);
\fill[udblue!4]  (8.9,-0.35)  rectangle (11.5,-10.2);
\fill[udblue!4]  (12.3,-0.35) rectangle (14.9,-10.2);

% --- Step 1 ---
\node[swimlane] (s1) at (0,-1.2) {Witnesses\\incident};
\node[swimlane] (s2) at (3.4,-1.2) {Opens app; selects\\incident type + location};
\draw[arr] (s1) -- (s2);

% --- Step 2 ---
\node[tech] (s3) at (6.8,-1.2) {Incident Service\\validates \& persists};
\draw[arr] (s2) -- (s3);

% --- Step 3: AI ---
\node[tech] (s4) at (6.8,-2.4) {AI plausibility\\scorer runs};
\draw[arr] (s3) -- (s4);

% --- Step 4: condition ---
\node[cond] (c1) at (6.8,-3.5) {Score\\OK?};
\draw[arr] (s4) -- (c1);

% --- Yes path ---
\node[tech] (s5) at (6.8,-4.8) {Crowd confirmation\\collected ($\geq$2)};
\draw[arr] (c1) -- node[right,font=\scriptsize]{Yes} (s5);

% --- No path ---
\node[swimlane] (s5b) at (10.2,-3.5) {Guard reviews\\low-score report};
\draw[arr] (c1.east) -- node[above,font=\scriptsize]{No/uncertain} (s5b);
\node[cond] (c2) at (10.2,-4.8) {Approve?};
\draw[arr] (s5b) -- (c2);
\draw[arr,dashed] (c2.west) -- node[above,font=\scriptsize]{No: reject} (6.8,-4.8);
\draw[arr] (c2.south) -- node[right,font=\scriptsize]{Yes} (10.2,-5.8);

% --- Step 5: dispatch ---
\node[tech] (s6) at (6.8,-6.0) {Priority score computed\\(zone $\times$ time $\times$ severity)};
\draw[arr] (s5) -- (s6);

\node[tech] (s7) at (6.8,-7.2) {Dispatcher broadcasts\\geofenced alert};
\draw[arr] (s6) -- (s7);

% --- Guard receives ---
\node[swimlane] (s8) at (10.2,-7.2) {Receives alert;\\deploys response};
\draw[arr] (s7.east) -- (s8.west);
\draw[arr] (10.2,-5.8) -- (s8);

% --- Community receives ---
\node[swimlane] (s9) at (3.4,-7.2) {Receives\\push alert};
\draw[arr] (s7.west) -- (s9.east);

% --- Escalation ---
\node[cond] (c3) at (10.2,-8.4) {Escalate?};
\draw[arr] (s8) -- (c3);
\node[swimlane] (s10) at (13.6,-8.4) {Authorities\\notified};
\draw[arr] (c3.east) -- node[above,font=\scriptsize]{Yes} (s10.west);

% --- Logging ---
\node[tech] (s11) at (6.8,-9.4) {Incident logged; analytics\\updated; reporter notified};
\draw[arr] (s7) -- (s11);
\draw[arr,dashed] (c3.south) |- (s11.east);

\end{tikzpicture}
\caption{CONOPS swim-lane diagram. Left lanes (blue background) represent
human actors; centre lane (gold background) represents the technology
platform. Dashed arrows indicate conditional flows.
Source: own elaboration (2026).}
\label{fig:conops}
\end{figure*}

\subsection{Emergency SOS workflow.}
When a user activates the SOS button (FR-08), the normal verification
pipeline is bypassed entirely. GPS coordinates are captured
automatically (FR-06), an immediate dispatch message is pushed to all
on-duty security guards, and a 90-second countdown begins. If no guard
acknowledges the alert within 90\,s, the system automatically escalates
to local authorities via the secure REST endpoint. This bypass design
is justified by the W1 finding that the current 11.4-minute latency is
incompatible with life-safety emergencies.

\subsection{Operational conditions and constraints.}
Three environmental constraints shape every operational decision:
(a)~\textbf{Night-time window} 18:00--22:00 has the highest observed
incident concentration (W1 direct observation: both suspicious events
in this window); the system must remain fully staffed and fully
responsive during this window (NFR-03).
(b)~\textbf{Connectivity dead zones}: the W1 observation phase did not
assess mobile signal coverage; offline queuing is therefore mandatory.
(c)~\textbf{Staff availability}: a 90-second moderator timeout prevents
the verification pipeline from stalling when no moderator is logged
in, escalating automatically to the on-call path.

%======================================================================
\section{System Architecture}
\label{sec:sysarch}

A systems architecture is not the same as a software architecture.
This section describes the complete system: stakeholders, their
interfaces with the technology, information flows between
organisational and technological components, and the relationships
between the system and its environment. The software architecture
(Section~\ref{sec:softarch}) specifies the technological components
in isolation.

\subsection{System context.}
Fig.~\ref{fig:syscontext} defines the system boundary. Everything
inside the boundary is managed, designed, and operated by the CSAS
project. Everything outside the boundary is the environment: external
actors, pre-existing platforms, and physical infrastructure.

\begin{figure}[h]
\centering
\begin{tikzpicture}[
  font=\footnotesize,every node/.style={align=center},
  ext/.style={rounded corners,draw=udblue!70,fill=udblue!8,
              inner sep=4pt,minimum width=2.0cm,minimum height=0.7cm},
  sys/.style={rounded corners,draw=udgold,thick,fill=udgold!10,
              inner sep=4pt,minimum width=2.0cm,minimum height=0.7cm},
  intf/.style={-{Latex[length=2mm]},thick},
  bintf/.style={Latex[length=2mm]-{Latex[length=2mm]},thick}
]
% System boundary
\draw[udgold,thick,dashed,rounded corners=8pt] (-2.2,-2.4) rectangle (2.2,2.4);
\node[anchor=north west,udgold,font=\small\itshape] at (-2.2,2.4) {System boundary};

% Internal nodes
\node[sys] (api)  at (0,1.2)  {API Gateway};
\node[sys] (svc)  at (0,0)    {Microservices\\(5 services)};
\node[sys] (db)   at (0,-1.2) {Data Layer\\(DB + Cache)};
\draw[intf] (api) -- (svc);
\draw[intf] (svc) -- (db);

% External actors
\node[ext] (users)  at (-5.5,1.5)  {Community\\Members};
\node[ext] (guard)  at (-5.5,0)    {Security\\Staff};
\node[ext] (mgmt)   at (-5.5,-1.5) {University\\Management};
\node[ext] (siure)  at (5.5,1.5)   {SIURE UD\\(existing)};
\node[ext] (auth)   at (5.5,0)     {Local Authorities\\(Police / EMS)};
\node[ext] (sms)    at (5.5,-1.5)  {SMS Gateway\\(Twilio)};

% Interfaces
\draw[bintf] (users.east) -- node[above,font=\scriptsize]{Reports / Alerts}   (-2.2,1.5);
\draw[bintf] (guard.east) -- node[above,font=\scriptsize]{Reviews / Actions}  (-2.2,0);
\draw[intf]  (api.west)   -- node[below,font=\scriptsize]{Analytics}          (mgmt.east);
\draw[bintf] (api.east)   -- node[above,font=\scriptsize]{Incidents / History}(siure.west);
\draw[intf]  (api.east)   -- node[above,font=\scriptsize]{Escalated alerts}   (auth.west);
\draw[intf]  (api.east)   -- node[above,font=\scriptsize]{Push fallback}      (sms.west);

\end{tikzpicture}
\caption{System context diagram. Blue boxes are external actors.
Gold boxes are system-managed components.
Bidirectional arrows indicate information flows in both directions.
Source: own elaboration (2026).}
\label{fig:syscontext}
\end{figure}

\subsection{Information flows between organisational and technological components.}
Table~\ref{tab:infoflow} maps each information flow to the
organisational sender, the system component that handles it, and the
downstream organisational receiver. This table answers the professor's
requirement to show how human processes interact with the technological
platform.

\begin{table}[h]
\caption{Organisational--Technology Information Flows}
\label{tab:infoflow}
\centering
\begin{tabularx}{\textwidth}{lXXX}
\toprule
\textbf{Flow} & \textbf{Org.\ sender} & \textbf{System component} & \textbf{Org.\ receiver} \\
\midrule
Incident report & Community member & Incident Service → Verification Engine & Security guard (alert) \\
SOS activation & Community member & Incident Service (bypass) & Security guard + Authorities \\
Report confirmation & Community member & Verification Engine (crowd logic) & Security guard (decision support) \\
Alert broadcast & System (Dispatcher) & Push / SMS engine & Community members in 500\,m \\
Incident escalation & Security supervisor & Integration layer → REST endpoint & Local authorities (Police / EMS) \\
Historical feed & SIURE UD (existing) & Integration layer → Analytics & Security supervisor, Mgmt \\
Heatmap report & System (Analytics) & Analytics engine → Dashboard & University management \\
Audit log & System & Audit service & System administrator \\
\bottomrule
\end{tabularx}
\end{table}

\subsection{Inputs, outputs, and constraints.}
Table~\ref{tab:ioc} specifies the system-level inputs, outputs, and
constraints, required for any compliant systems engineering description.

\begin{table}[h]
\caption{System Inputs, Outputs, and Operational Constraints}
\label{tab:ioc}
\centering
\begin{tabularx}{\textwidth}{lX}
\toprule
\textbf{Category} & \textbf{Items} \\
\midrule
\textbf{Inputs} & User registration data; incident reports (type, GPS, timestamp, media); community confirmation signals; SOS activations; SIURE UD historical feed; weather API data; guard acknowledgement signals \\
\textbf{Outputs} & Geofenced push alerts; SMS fallback alerts; administrative dashboard events; heatmap reports; incident audit log entries; escalation messages to authorities \\
\textbf{Operational constraints} & Ley 1581/2012 data protection; TLS 1.3 encryption; 99.9\% uptime SLA; $\leq$3-interaction reporting (NFR-06); maximum 45-second verification budget; 90-second escalation timeout \\
\textbf{Environmental constraints} & Unreliable mobile connectivity in parts of campus (assumption unvalidated, see Section~\ref{sec:assumptions}); limited security staff availability outside business hours; no existing standardised incident classification schema \\
\bottomrule
\end{tabularx}
\end{table}

%======================================================================
\section{Software Architecture}
\label{sec:softarch}

\subsection{Architectural style.}
The CSAS adopts a \textbf{microservices, event-driven architecture}.
This choice is defended in DDR-03 (Section~\ref{sec:ddr}); briefly,
the verification, notification, and analytics subsystems have
independent load profiles and scaling needs that make a monolithic
deployment suboptimal, and the alert dispatch must be decoupled from
incident ingestion to prevent cascading delays during high-volume
periods.

\subsection{Five-layer organisation.}
The architecture is organised across five layers:
(i)~\textbf{Presentation} --- mobile app (React Native) and web
portal;
(ii)~\textbf{Application} --- API gateway (Node.js / Nginx) for
authentication, rate limiting, and service routing;
(iii)~\textbf{Service} --- five domain microservices
(Table~\ref{tab:services});
(iv)~\textbf{Data} --- PostgreSQL + PostGIS database, Redis cache,
and media object store;
(v)~\textbf{Integration} --- bridges to SIURE UD, local authorities
REST endpoint, and Twilio SMS gateway.

\begin{table}[h]
\caption{Microservices --- Responsibilities and Domain Events}
\label{tab:services}
\centering
\begin{tabularx}{\textwidth}{lXl}
\toprule
\textbf{Service} & \textbf{Responsibility} & \textbf{Domain events} \\
\midrule
Incident & Receives, validates, and persists reports; enriches with geospatial metadata. & Publishes \texttt{IncidentCreated} \\
Verification & Orchestrates the hybrid AI + human pipeline (DDR-01). Adaptive threshold, priority bypass for SOS and high-severity types. & Consumes \texttt{IncidentCreated}; publishes \texttt{IncidentVerified} \\
Dispatcher & Computes priority score (DDR-02); routes to push / SMS. & Consumes \texttt{IncidentVerified}; publishes \texttt{AlertDispatched} \\
User & Manages registration, roles, alert preferences; implements alert-fatigue dampening. & Consumes \texttt{AlertDispatched} \\
Analytics & Aggregates domain events into the data warehouse; produces heatmaps. & Consumes all events \\
\bottomrule
\end{tabularx}
\end{table}

\subsection{Technology stack.}
Table~\ref{tab:stack} presents the recommended technology stack.
Every choice cites the requirement it satisfies; the deeper
engineering justification is in the DDRs.

\begin{table}[h]
\caption{Recommended Technology Stack}
\label{tab:stack}
\centering
\begin{tabularx}{\textwidth}{llX}
\toprule
\textbf{Component} & \textbf{Technology} & \textbf{Justification} \\
\midrule
Mobile app & React Native & Single codebase Android/iOS; $\leq$3-interaction reporting (NFR-06). \\
API gateway & Node.js / Nginx & High I/O throughput; handles concurrent submissions and rate limiting. \\
Backend services & Node.js & Asynchronous, event-driven I/O supports $<$30\,s latency (NFR-01). \\
Verification Engine & Python / FastAPI & Python ML ecosystem for AI plausibility scorer (DDR-01). \\
Message broker & RabbitMQ & Priority-queue support; decouples ingestion from verification (mitigates volume spike risk). \\
Database & PostgreSQL + PostGIS & Native geospatial queries for zone-weighted routing and heatmaps. \\
Push notifications & Firebase FCM & Cross-platform; supports 99.9\% delivery (NFR-02). \\
SMS fallback & Twilio API & Guaranteed delivery for users without push capability; equity of access. \\
Infrastructure & AWS / Azure + Kubernetes & Container orchestration enables auto-scaling (NFR-03, NFR-05). \\
\bottomrule
\end{tabularx}
\end{table}

\subsection{External integrations.}
Three external integrations extend the system's reach beyond the
campus platform:
\textbf{SIURE UD} (bidirectional REST bridge) provides the historical
baseline and complements the existing institutional workflow without
replacing it;
\textbf{Local authorities} (outbound secure REST endpoint) activated
only for high-severity escalations;
\textbf{SMS gateway} (Twilio) functions as both a delivery fallback
and an alternative submission channel for users without smartphone
access.

%======================================================================
\section{Design Decision Records}
\label{sec:ddr}

A design decision record (DDR) documents \emph{why} a specific
architectural choice was made: what alternatives were considered,
what criteria drove the evaluation, and under what conditions the
chosen approach could fail. This section provides the engineering
justification that was absent from the first submission.

%----------------------------------------------------------------------
\subsection{DDR-01 --- Hybrid AI + Human Verification Pipeline.}
\label{sec:ddr01}

\textbf{Decision:} Implement a two-stage verification pipeline
combining an AI plausibility scorer (Stage~1) with human moderator
review (Stage~2).

\textbf{Evidence driving the decision:} W1 sensitivity analysis
identified that ``if too many alerts are sent, users may start
ignoring them'' (Balancing Loop B1). Any single-stage mechanism
that is either too permissive (false positives trigger B1) or too
strict (genuine reports suppressed, Loop R1 stalls) creates a failure
mode.

\textbf{Alternatives considered:}
Table~\ref{tab:ddr01} evaluates four options on five criteria.

\begin{table}[h]
\caption{DDR-01 --- Evaluation of Verification Alternatives}
\label{tab:ddr01}
\centering
\begin{tabularx}{\textwidth}{lXXXXX}
\toprule
\textbf{Option} & \textbf{Latency} & \textbf{Accuracy} & \textbf{Failure resilience} & \textbf{Alert-fatigue risk} & \textbf{Cost} \\
\midrule
AI-only & Very low & Medium (degrades silently) & Low (no fallback) & High if model errors & Low \\
Human-only & High ($>$2 min estimated) & High & High & Low & High (staffing) \\
Crowd-only ($\geq$3 confirmations) & Variable & Medium & Low (sparse off-hours) & Medium & Very low \\
\textbf{Hybrid AI + human (chosen)} & Low (AI pre-filters) & High (human corrects AI) & High (human override) & Low (AI suppresses noise) & Medium \\
\bottomrule
\end{tabularx}
\end{table}

\textbf{Selection rationale:} AI-only has no accountability and
silent failure modes. Human-only cannot meet the $<$30\,s latency
requirement (NFR-01) during off-hours when staff availability is low.
Crowd-only degrades during nights and early mornings precisely when
the W1 data shows the highest risk (all suspicious events were during
evening/night). Hybrid combines the speed of AI pre-filtering with
the authority and accuracy of human decision-making for uncertain
cases.

\textbf{Conditions of failure:}
(1)~AI model quality degrades due to dataset drift; mitigated by
monthly retraining on verified CSAS events.
(2)~No moderator available during the 45-second verification budget;
mitigated by an automatic timeout that escalates to the on-call
security guard.
(3)~Community confirmations arrive slowly for low-population
off-hours; mitigated by the staff override path that bypasses the
crowd requirement.

\textbf{How effectiveness will be measured:}
False-positive rate on verified alerts (target: $<$5\%);
average verification latency (target: $<$45\,s);
daily human-override rate as a proxy for AI quality.

%----------------------------------------------------------------------
\subsection{DDR-02 --- Zone-Weighted and Time-Weighted Priority Dispatch.}
\label{sec:ddr02}

\textbf{Decision:} Compute a priority score
$P = S_{\text{base}} \times W_z \times W_t$
where $S_{\text{base}}$ is the incident severity class,
$W_z$ is the zone weight, and $W_t$ is the time-of-day weight,
before deciding which queue tier and notification radius to use.

\textbf{Evidence driving the decision:}
W1 identified that ``alert activity is not uniformly distributed
across the campus or the day'' and that ``the system should apply
zone-weighted and time-weighted priority logic rather than treating
all incoming reports identically.'' Specifically:
the parking lot was the highest-risk zone across all observation
sessions; both suspicious events occurred at night; both coincided
with rainy conditions.

\textbf{Model definition:}
Table~\ref{tab:weights} specifies the calibrated weight values.
Initial values are derived from W1 observations; they will be revised
quarterly using incident analytics (FR-09).

\begin{table}[h]
\caption{DDR-02 --- Priority Score Weight Tables}
\label{tab:weights}
\centering
\begin{tabular}{llc}
\toprule
\multicolumn{3}{l}{\textbf{Severity base score $S_{\text{base}}$}} \\
\midrule
SOS / life-safety & & 4.0 \\
Theft / assault & & 2.0 \\
Harassment & & 1.5 \\
Suspicious activity & & 1.0 \\
\midrule
\multicolumn{3}{l}{\textbf{Zone weight $W_z$ (W1 observation data)}} \\
\midrule
Parking lot (night) & Highest-risk across all sessions & 1.5 \\
Dormitory area & Poor lighting; latent risk & 1.3 \\
Main entrance & High foot traffic; deterrent effect & 1.0 \\
Library area & High activity; no suspicious events & 1.0 \\
\midrule
\multicolumn{3}{l}{\textbf{Time weight $W_t$ (W1 observation + survey)}} \\
\midrule
18:00--22:00 & All suspicious events in this window & 1.5 \\
22:00--06:00 & Very low foot traffic; high risk & 1.3 \\
06:00--18:00 & Normal activity; lower risk & 1.0 \\
\bottomrule
\end{tabular}
\quad
\begin{tabular}{lll}
\toprule
\multicolumn{3}{l}{\textbf{Priority tier and dispatch action}} \\
\midrule
$P \geq 4.0$ & CRITICAL & Immediate dispatch; 500\,m alert; auth. escalation path \\
$2.5 \leq P < 4.0$ & HIGH & Standard dispatch; 500\,m alert; on-duty guard \\
$1.0 \leq P < 2.5$ & STANDARD & Queued dispatch; 250\,m alert \\
$P < 1.0$ & SUPPRESSED & Logged; not broadcast \\
\bottomrule
\end{tabular}
\end{table}

\textbf{Alternatives considered:}
(1)~Flat 500\,m broadcast regardless of priority --- produces alert
storms in high-risk zones, accelerating B1 (alert fatigue).
(2)~Distance-only routing with no time or zone weights --- ignores the
strong temporal and spatial patterns documented in W1.
(3)~Manual categorisation by guards --- too slow for NFR-01 ($<$30\,s)
and depends on guards being logged in.
The multiplicative model was chosen because it is transparent,
calibratable from real data, and can be updated periodically
without redeployment (weights stored in a configuration table).

\textbf{Conditions of failure:}
(1)~Stale weights that no longer reflect current risk distribution;
mitigated by quarterly recalibration.
(2)~GPS inaccuracy places the incident in the wrong zone;
mitigated by a 20\,m tolerance buffer at zone boundaries.
(3)~Extreme simultaneous volume (many high-P reports at once)
saturates the CRITICAL queue; mitigated by RabbitMQ priority lanes
and auto-scaling (see DDR-03).

\textbf{How effectiveness will be measured:}
Average alert latency by priority tier (target: CRITICAL $<$15\,s);
false-positive rate by tier; guard response time after alert receipt.

%----------------------------------------------------------------------
\subsection{DDR-03 --- Microservices vs.\ Alternative Architectural Styles.}
\label{sec:ddr03}

\textbf{Decision:} Decompose the system into five independently
deployable microservices communicating over a RabbitMQ message broker.

\textbf{Alternatives considered:}
Table~\ref{tab:ddr03} compares the four candidate styles.

\begin{table}[h]
\caption{DDR-03 --- Evaluation of Architectural Style Alternatives}
\label{tab:ddr03}
\centering
\begin{tabularx}{\textwidth}{lXXXX}
\toprule
\textbf{Style} & \textbf{Independent scaling} & \textbf{Deployment complexity} & \textbf{Failure isolation} & \textbf{Team fit} \\
\midrule
Monolith & None (all-or-nothing) & Low & Low (one failure cascades) & Low (parallel dev difficult) \\
Serverless (FaaS) & High & Medium & High & Medium (cold-start latency risks NFR-01) \\
SOA (heavyweight) & Medium & High & Medium & Low (requires ESB) \\
\textbf{Microservices (chosen)} & Per-service & Medium & High & High (teams own services) \\
\bottomrule
\end{tabularx}
\end{table}

\textbf{Selection rationale:}
Verification and notification have non-correlated load spikes:
verification peaks when many users report simultaneously; notification
peaks when a major alert is dispatched. A monolith cannot scale these
independently. Serverless introduces cold-start latency that could
violate NFR-01 for CRITICAL alerts. SOA imposes integration overhead
with no benefit over microservices at this scale.

\textbf{Conditions of failure:}
Network partition between services could delay the
\texttt{IncidentCreated} → Verification event; mitigated by RabbitMQ
durability guarantees and a dead-letter queue that re-enqueues
unprocessed events.

%======================================================================
\section{Operational Scenarios}
\label{sec:scenarios}

The priority scoring model and verification pipeline must be
grounded in concrete operational reasoning. The following three
scenarios demonstrate the system's decision logic with specific inputs,
calculated outcomes, and failure modes. All scenarios use the weight
tables in DDR-02 (Table~\ref{tab:weights}).

%----------------------------------------------------------------------
\subsection{Scenario 1 --- Evening parking-lot theft.}
\textbf{Context:} 21:30, rainy night. A student at the parking lot
witnesses a theft and submits a report with GPS coordinates placing
the incident in the Parking Lot zone.

\textbf{Priority calculation:}
\[
P = S_{\text{base}} \times W_z \times W_t = 2.0 \times 1.5 \times 1.5 = \mathbf{4.5 \rightarrow \text{CRITICAL}}
\]

\textbf{System response:}
The Incident Service persists the report and publishes
\texttt{IncidentCreated}. The Verification Engine routes it to the
priority bypass path (CRITICAL tier skips the 45-second crowd
confirmation requirement). The AI scorer runs in parallel ($<$3\,s)
and returns a plausibility score above threshold. The Dispatcher
broadcasts a push alert to all users within 500\,m and
simultaneously sends an alert to the on-duty guard's device.
\emph{Expected latency: 8--12\,s from report to first notification.}

\textbf{Fallback:} If rain has caused the user's GPS signal to degrade
(low-accuracy reading), the 20\,m zone buffer places the incident in
the correct zone. If push delivery fails, the SMS fallback activates
after 10\,s.

%----------------------------------------------------------------------
\subsection{Scenario 2 --- SOS emergency activation.}
\textbf{Context:} 20:45. A student near the dormitory area activates
the SOS button (FR-08). No incident description is provided.

\textbf{Priority calculation:}
\[
P = 4.0 \times 1.3 \times 1.5 = \mathbf{7.8 \rightarrow \text{CRITICAL (bypass)}}
\]

\textbf{System response:}
The Incident Service captures GPS coordinates automatically and
publishes a \texttt{SOSActivated} domain event. The Verification
Engine's SOS handler bypasses all confirmation logic and routes
directly to the Dispatcher within $<$2\,s. The Dispatcher broadcasts
a CRITICAL alert to the entire campus and simultaneously notifies
all on-duty security guards. A 90-second ACK timer starts. If no
guard acknowledges within 90\,s, the integration layer pushes an
escalation request to the local-authorities REST endpoint.
\emph{Expected latency: $<$4\,s from SOS press to first guard notification.}

%----------------------------------------------------------------------
\subsection{Scenario 3 --- High-volume surge (post-event crowd).}
\textbf{Context:} 22:15 following a large campus event. Thirty users
simultaneously submit reports related to a physical altercation near
the main entrance.

\textbf{Priority calculations (representative sample):}
\begin{itemize}
\item Assault reports at main entrance at 22:15:
  $P = 2.0 \times 1.0 \times 1.3 = 2.6$ → HIGH
\item Suspicious-activity reports (noise from crowd):
  $P = 1.0 \times 1.0 \times 1.3 = 1.3$ → STANDARD
\end{itemize}

\textbf{System response:}
RabbitMQ routes the 30 reports into two priority lanes: HIGH (assault
reports) processed first, STANDARD queued behind. The spatial
clustering algorithm groups geographically co-located STANDARD reports
within a 30\,m radius into a single composite alert, reducing 22
individual notifications to 4 composite ones. Kubernetes detects
queue depth exceeding 15 events and triggers a scale-out policy on
the Verification Engine, adding two pod replicas within 45\,s.
\emph{Expected outcome: the 8 HIGH-priority reports are dispatched in
$<$30\,s (NFR-01 met); composite STANDARD alerts dispatched in
$<$90\,s.}

\textbf{Failure mode under extreme surge:} If 300+ simultaneous
reports arrive (stress scenario, NFR-05), the CRITICAL/HIGH queues
remain unaffected by priority isolation; STANDARD alerts may
experience latency up to 5 minutes but are never lost (durable queue).
Load-test targets for Phase~3 validate this envelope.

%======================================================================
\section{Risk and Complexity Management}
\label{sec:risk}

\subsection{Sensitivity parameter mitigation.}
Table~\ref{tab:risk} maps each of the five sensitivity parameters
identified in W1 to its quantified risk and its design-level mitigation
mechanism, showing the traceability from analysis to design.

\begin{table}[h]
\caption{Sensitivity Parameter Risk and Design Mitigation}
\label{tab:risk}
\centering
\begin{tabularx}{\textwidth}{lXX}
\toprule
\textbf{Parameter} & \textbf{Risk} & \textbf{Design mechanism} \\
\midrule
Report volume spike & Verification engine saturates; $>$30\,s latency. & RabbitMQ priority lanes; auto-scaling on Verification pods; SOS and CRITICAL bypass standard queue. \\
Low user participation & Real incidents undetected; R1 loop stalls. & $\leq$3-interaction interface (FR-01); anonymous reporting option; community engagement module; SIURE integration as institutional endorsement. \\
Night-time concentration & Higher alert volume; fewer moderators online. & Time weights elevate night-hour priority; on-call escalation path; 45-second moderator timeout with auto-escalation. \\
High-risk zone clustering & Alert storms in parking lot and dormitory. & Zone-based rate limiting; spatial clustering aggregates co-located simultaneous reports (max 1 composite per 30\,m / 60\,s). \\
Adverse weather & Visibility drops; GPS degrades; incident risk rises. & Weather API feeds risk scoring model; adverse conditions elevate $W_t$ by one tier in affected zones; GPS 20\,m boundary tolerance. \\
\bottomrule
\end{tabularx}
\end{table}

\subsection{Failure mode analysis.}
Table~\ref{tab:failures} maps each major architectural dependency to
its failure mode, impact, and mitigation.

\begin{table}[h]
\caption{Failure Mode Analysis}
\label{tab:failures}
\centering
\begin{tabularx}{\textwidth}{lXXX}
\toprule
\textbf{Failure mode} & \textbf{Impact} & \textbf{Mitigation} & \textbf{Recovery time target} \\
\midrule
API gateway outage & All submissions fail & Multi-region deployment + health checks; mobile app queues reports locally. & $<$60\,s automatic failover \\
Verification Engine overload & Alert latency exceeds NFR-01 & Auto-scaling pods; CRITICAL/SOS bypass path is unaffected. & $<$45\,s pod spin-up \\
Push delivery failure & Users miss critical alerts & SMS fallback activates after 10\,s push timeout. & $<$10\,s \\
Database unavailability & Incident data loss risk & Primary-replica replication; write-ahead logging; daily backups. & Zero data loss; $<$120\,s reconnect \\
AI scorer false positive & Legitimate incident blocked & Human review queue; 90-second escalation timer if no moderator. & $<$90\,s override \\
Mobile dead zone & Submission fails & Offline queue in app; SMS submission channel as fallback. & On reconnect \\
\bottomrule
\end{tabularx}
\end{table}

%======================================================================
\section{Performance Targets and Quality Assurance}

\subsection{Key Performance Indicators.}
Table~\ref{tab:kpi} defines the measurable targets. Each KPI is
aligned with a specific requirement and with a W1 operational gap.

\begin{table}[h]
\caption{Key Performance Indicators}
\label{tab:kpi}
\centering
\begin{tabularx}{\textwidth}{lllXl}
\toprule
\textbf{KPI} & \textbf{Target} & \textbf{Req.} & \textbf{W1 gap} & \textbf{Freq.} \\
\midrule
Alert latency (end-to-end) & $<$30\,s P95 & NFR-01 & Response time $-$6.4 min & Daily \\
Delivery success rate & $\geq$99.9\% & NFR-02 & Delivery $-$12.9\,pp & Daily \\
System availability & $\geq$99.9\% & NFR-03 & --- & Monthly \\
Incident record rate & $\geq$85\% logged & FR-07 & Completeness $+$44\,pp & Weekly \\
User adoption & $\geq$60\% students & --- & Adoption $+$48\,pp & Monthly \\
Verification latency & $<$45\,s & DDR-01 & --- & Daily \\
False-positive rate & $<$5\% verified alerts & DDR-01 & B1 loop mitigation & Daily \\
Scalability (surge) & 300\% surge, 0 timeouts & NFR-05 & --- & Monthly load test \\
\bottomrule
\end{tabularx}
\end{table}

\subsection{Testing strategy.}
\textbf{Unit testing} validates each microservice in isolation;
minimum 80\% code coverage per service.
\textbf{Integration testing} verifies the full event chain:
Incident~Service → Verification~Engine → Dispatcher → Push/SMS,
running against a staging environment seeded with the four observed
campus zones.
\textbf{Load testing} simulates a 300\% concurrent-request surge
(NFR-05) with target: zero timeout errors and $<$30\,s latency at P95.
A secondary test simulates the Scenario~3 volume spike (30
simultaneous reports) to validate RabbitMQ priority isolation.
\textbf{Operational scenario testing} replays the three scenarios in
Section~\ref{sec:scenarios} against the prototype, verifying that
priority scores, dispatch paths, and latency targets match the
specified values.
\textbf{User acceptance testing (UAT)} uses a pilot cohort of students
and security staff; success criterion: $\geq$85\% of participants
submit a report in $<$2 minutes and rate usability $\geq$4/5.
\textbf{Security testing} verifies TLS 1.3 compliance, Ley~1581
data controls, access-controlled audit log, and anonymous-reporter
encryption.

%======================================================================
\section{Implementation Roadmap}
\label{sec:roadmap}

The implementation is organised in four phases. Each phase specifies
deliverables, success criteria, and the open assumptions
(Section~\ref{sec:assumptions}) that must be validated before the
next phase begins.

\textbf{Phase 1 --- Planning and architecture (Month 1).}
Deliverables: finalised architecture diagrams; complete requirements
traceability matrix; development environment and CI/CD pipeline setup;
campus-wide mobile connectivity audit.
Success criterion: every FR and NFR traceable to a W1 finding;
connectivity dead-zones mapped and offline-queuing strategy confirmed.
W1 assumption to close: A3 (mobile connectivity, unvalidated).

\textbf{Phase 2 --- Core development (Month 2).}
Deliverables: functional prototype with FR-01 (incident reporting)
and FR-02 (geofenced alert) operational end-to-end;
database schema; message broker configuration; all five microservices
in development.
Success criterion: a student can submit a theft report and receive a
geofenced push alert within 30\,s on the staging environment.

\textbf{Phase 3 --- Testing and UAT (Month 3).}
Deliverables: unit, integration, load, and operational-scenario test
reports; UAT report from pilot cohort; community engagement materials.
Success criterion: NFR-01 ($<$30\,s) and NFR-02 (99.9\%) met under
simulated 300\% surge; $\geq$85\% UAT success rate.
W1 assumption to close: A1 (willingness to report under real
conditions).

\textbf{Phase 4 --- Deployment and monitoring (Month 4).}
Deliverables: production deployment on cloud infrastructure;
SIURE UD integration live; real-time monitoring dashboard active;
Ley~1581 data-processing agreement signed.
Success criterion: 99.9\% uptime over 30-day monitoring window;
guard acknowledgement rate $\geq$90\% on CRITICAL alerts.
W1 assumption to close: A2 (verification latency within 45-second
budget under real load) and A5 (institutional integration).

%======================================================================
\section{Ethical and Safety Considerations}

\subsection{Privacy and data protection.}
The CSAS processes three categories of personal data: registration
information (name, institutional email, campus role); real-time
location data during incident submission; and incident media (photos,
video). Each category carries distinct obligations under Ley~1581 de
2012, which is directly binding on the institution.

Four compliance requirements are enforced by design: explicit informed
consent at registration (purpose-specific); data used only for
declared safety purposes, no third-party transfer or advertising use;
user rights of access, correction, and deletion supported by the User
Service; and a 12-month incident retention policy with irreversible
anonymisation thereafter.

For sensitive incident types (harassment, threats), the anonymous
reporting option encrypts reporter identity at rest, accessible only
to system administrators under a documented, access-controlled
protocol with immutable audit-log entries.

\subsection{Equity and accessibility.}
A mobile-first design must not create an equity gap. The SMS gateway
(Twilio) functions as both a delivery fallback and a submission
channel, ensuring that community members without smartphone access
can both report incidents and receive alerts. The web portal provides
a WiFi-connected alternative. All user-facing components use Spanish
as the primary language; the onboarding flow must be completable in
under two minutes (NFR-06) without presupposing advanced digital
literacy.

\subsection{Prevention of misuse.}
Three misuse vectors are explicitly addressed:
\textbf{False reports:} mitigated by the hybrid verification pipeline
and human moderation (DDR-01).
\textbf{Harassment through the platform:} deterred by the immutable
audit log that records every moderation action.
\textbf{Anonymous report flooding:} rate limiting on the anonymous
path (maximum 3 anonymous reports per user per hour) prevents a
single actor from overwhelming the verification queue.

\subsection{Unintended consequences.}
Four classes of unintended consequences, identified in the W1 revised
submission, are actively mitigated:
\textbf{Vigilantism:} alert content uses advisory language
(``Incident reported nearby; please take precautions'') and never
directive language (``Approach and investigate'').
\textbf{Risk profiling:} the analytics layer enforces aggregation
policies; no individual movement data is exposed in heatmaps.
\textbf{Reduced personal vigilance:} onboarding explicitly frames
the CSAS as a complement to, not a substitute for, personal
situational awareness.
\textbf{Fear amplification:} the first operational period includes
a communication campaign contextualising alert frequency relative
to campus-wide incident prevalence.

%======================================================================
\section{Open Assumptions Requiring Validation}
\label{sec:assumptions}

Table~\ref{tab:assumptions} restates the five assumptions from the
W1 revised submission, updated with the specific validation action
assigned to each implementation phase.

\begin{table}[h]
\caption{Open Assumptions and Validation Plan}
\label{tab:assumptions}
\centering
\begin{tabularx}{\textwidth}{clXlX}
\toprule
\textbf{\#} & \textbf{Assumption} & \textbf{Current evidence} & \textbf{Status} & \textbf{Validation action} \\
\midrule
A1 & Users will report if friction is reduced & 100\% of incident-experienced survey respondents willing to adopt (W1, $n=20$) & Plausible & Phase 3 UAT: measure actual submission rate in pilot \\
A2 & Verification can operate within 45\,s budget & Scenario analysis; no empirical measurement & Untested & Phase 4 load test under real conditions \\
A3 & Mobile connectivity reliable across all zones & Not assessed in W1 & Unvalidated & Phase 1: connectivity audit before deployment \\
A4 & False-positive rate below B1 threshold & Theoretical; DDR-01 design provides mitigation & Untested & Phase 3: measure FP rate in staging \\
A5 & Institutional integration (SIURE, security team) & Staff interviews supportive; no formal commitment & Partially supported & Phase 1: data-processing agreement + admin onboarding \\
\bottomrule
\end{tabularx}
\end{table}

%======================================================================
\section{Conclusions}

This document has translated the empirical findings of Workshop~No.~1
into a complete, traceable systems design specification for the CSAS.
Five conclusions summarise the design-phase contributions.

\textbf{The design is system-driven, not feature-driven.} Every
functional requirement is traceable to a W1 empirical finding
(Table~\ref{tab:fr}); every architectural choice is justified by a
Design Decision Record (Section~\ref{sec:ddr}) that documents
alternatives considered and conditions of failure.

\textbf{The priority dispatch model is now formally specified.} The
multiplicative model $P = S_{\text{base}} \times W_z \times W_t$
converts the W1 observation-based risk findings into an operational
mechanism with explicit weights, thresholds, and calibration policy.
The model is transparent, updatable, and measurable.

\textbf{The verification pipeline resolves the B1--R1 tension.} The
hybrid AI + human design (DDR-01) is the only configuration that
simultaneously achieves low latency (NFR-01), high accuracy, and
resilience to moderator unavailability. This conclusion is supported
by the alternative evaluation in Table~\ref{tab:ddr01}.

\textbf{Three concrete scenarios confirm the design's operational
validity.} The scenarios in Section~\ref{sec:scenarios} demonstrate
that the system produces correct, calibrated responses to the three
most important risk events identified in W1: an evening parking-lot
incident, a life-safety SOS, and a high-volume post-event surge.

\textbf{The largest remaining risks are behavioural, not technical.}
The adoption gap ($+$48\,pp to reach the 60\% target) cannot be
closed by technical means alone; it requires institutional endorsement,
community engagement, and a sustained communication strategy. This
is the primary operational risk for the next implementation phase.

%======================================================================
\section*{References}
\addcontentsline{toc}{section}{References}
\begin{enumerate}[label={[\arabic*]},leftmargin=1.5em,itemsep=2pt]
\item C.\,A.\ Sierra, \emph{TeamWork as Computer Engineers --- CS Collaboration Guidelines}. Bogotá: Universidad Distrital Francisco José de Caldas, 2026.
\item P.\ Checkland, \emph{Systems Thinking, Systems Practice}. New York: John Wiley \& Sons, 1981.
\item D.\,H.\ Meadows, \emph{Thinking in Systems: A Primer}. White River Junction, VT: Chelsea Green Publishing, 2008.
\item ISO/IEC/IEEE, \emph{Systems and Software Engineering --- System Life Cycle Processes}, Std.\ 15288:2015, 2015.
\item J.\,W.\ Creswell and J.\,D.\ Creswell, \emph{Research Design: Qualitative, Quantitative, and Mixed Methods Approaches}, 5th ed.\ Thousand Oaks, CA: SAGE Publications, 2018.
\item Congreso de Colombia, \emph{Ley 1581 de 2012 --- Régimen General de Protección de Datos Personales}, 2012.
\item L.\ Bass, P.\ Clements, and R.\ Kazman, \emph{Software Architecture in Practice}, 3rd ed.\ Boston: Addison-Wesley, 2013.
\item N.\ Dragoni et al., ``Microservices: Yesterday, Today, and Tomorrow,'' \emph{Present and Ulterior Software Engineering}, pp.\ 195--216, 2017.
\item J.\,D.\ Sterman, \emph{Business Dynamics: Systems Thinking and Modeling for a Complex World}. Boston: McGraw-Hill, 2000.
\end{enumerate}

\end{document}
